\documentclass[a4paper,11pt]{amsart}

\usepackage[margin=1in]{geometry}
\usepackage{amsmath,amssymb,amsthm,mathtools}
\usepackage{enumitem}
\usepackage[colorlinks=true,linkcolor=blue,citecolor=blue,urlcolor=blue]{hyperref}

\newtheorem{theorem}{Theorem}[section]
\newtheorem{proposition}[theorem]{Proposition}
\newtheorem{lemma}[theorem]{Lemma}
\newtheorem{corollary}[theorem]{Corollary}
\newtheorem{importedtheorem}[theorem]{Imported theorem}
\theoremstyle{remark}
\newtheorem{remark}[theorem]{Remark}

\newcommand{\Sub}{\operatorname{Sub}}
\newcommand{\Sym}{\operatorname{Sym}}
\newcommand{\F}{\mathbb F}
\newcommand{\Prob}{\mathbb P}
\newcommand{\gbinom}[2]{\genfrac{[}{]}{0pt}{}{#1}{#2}_2}

\title[Subgroups of the symmetric group]{Subgroups of the symmetric group: counting and order statistics}
\author{}
\date{}

\begin{document}

\begin{abstract}
Let
\[
        s(n)=|\{H:H\leq S_n\}|.
\]
The asymptotic formula is
\[
        \log_2 s(n)=\frac{n^2}{16}+o(n^2),
        \qquad\text{equivalently}\qquad
        s(n)=2^{(1/16+o(1))n^2}.
\]
The lower bound is elementary and comes from the subgroup lattice of an elementary abelian
2-subgroup of rank \(\lfloor n/2\rfloor\). The matching upper bound is the theorem of
Roney-Dougal and Tracey, which proves a strong form of Pyber's 1993 conjecture. We also record
what is presently available as a genuine global statistical theorem about the order of a uniformly
random subgroup \(H_n\leq S_n\): for every
\[
        \nu<\rho:=\frac12-\frac{\sqrt3}{4},
\]
one has \(\log_2 |H_n|\geq \nu n\) with probability tending to \(1\). This follows from their theorem
that the Sylow \(2\)-subgroup of a random subgroup has order at least \(2^{\nu n}\) with high
probability. Separately, inside the elementary abelian lower-bound model, a uniformly random
subgroup of \(C_2^{\lfloor n/2\rfloor}\leq S_n\) has the sharper model law
\(2^{n/4+o(n)}\). Finally, a dihedral-block construction shows why the leading
\(2^{n^2/16+o(n^2)}\) enumeration alone does not determine an exact global concentration law for
\(|H_n|\).
\end{abstract}

\maketitle

\section{Statement of the answer}

All logarithms are base \(2\). Let
\[
        s(n)=|\Sub(S_n)|
\]
denote the number of subgroups of the symmetric group on \(n\) letters.

\begin{theorem}[Asymptotic subgroup count]\label{thm:main}
As \(n\to\infty\),
\[
        s(n)=2^{n^2/16+o(n^2)}.
\]
Equivalently,
\[
        \log_2 s(n)\sim \frac{n^2}{16}.
\]
\end{theorem}

The proof has two different parts. The lower bound is elementary and is proved in full below. The
upper bound is a deep finite-permutation-group theorem of Roney-Dougal and Tracey. Historically,
Pyber proved a weaker upper bound and conjectured the constant \(1/16\); the exact asymptotic
was proved by Roney-Dougal and Tracey.

For the statistical part, throughout the paper \(H_n\) denotes a uniformly random element of the
finite set \(\Sub(S_n)\). A \emph{section} of a finite group \(G\) means a quotient \(A/B\), where
\(B\trianglelefteq A\leq G\).

\begin{theorem}[Global order statistic currently supplied by the enumeration theory]\label{thm:global-order}
Let
\[
        \rho=\frac12-\frac{\sqrt3}{4}=0.066987\ldots .
\]
For every \(\varepsilon>0\),
\[
        \Prob\left(\log_2 |H_n|\geq (\rho-\varepsilon)n\right)\longrightarrow 1.
\]
Moreover, deterministically,
\[
        \log_2 |H_n|\leq \log_2(n!)\leq n\log_2 n.
\]
Thus a uniformly random subgroup of \(S_n\) has exponentially large order with high probability:
\[
        |H_n|\geq 2^{(\rho-o(1))n}
\]
in the usual lower-tail sense that every smaller linear exponent holds with probability tending to
\(1\).
\end{theorem}

\begin{remark}
Theorem~\ref{thm:global-order} is a real statistical theorem about uniformly random subgroups of
\(S_n\), but it is not an exact concentration theorem of the form \(\log_2|H_n|\sim c n\). Such an
exact global order law is a finer question than the leading enumeration
\(s(n)=2^{n^2/16+o(n^2)}\). The model theorem in Section~\ref{sec:model-order} and the
construction in Section~\ref{sec:dihedral} explain this distinction.
\end{remark}

\section{Gaussian binomial coefficients}

For integers \(0\leq d\leq r\), let
\[
        \gbinom{r}{d}
\]
denote the number of \(d\)-dimensional linear subspaces of \(\F_2^r\). These numbers are the
Gaussian binomial coefficients over \(\F_2\).

\begin{lemma}[Formula for \(\gbinom{r}{d}\)]\label{lem:gauss-formula}
For \(0\leq d\leq r\),
\[
        \gbinom{r}{d}
        =
        \frac{(2^r-1)(2^r-2)\cdots(2^r-2^{d-1})}
        {(2^d-1)(2^d-2)\cdots(2^d-2^{d-1})}.
\]
\end{lemma}

\begin{proof}
Count ordered bases of \(d\)-dimensional subspaces of \(\F_2^r\) in two ways. An ordered linearly
independent \(d\)-tuple in \(\F_2^r\) can be chosen in
\[
        (2^r-1)(2^r-2)\cdots(2^r-2^{d-1})
\]
ways: after \(i\) independent vectors have been chosen, their span has \(2^i\) elements, so the
next vector has \(2^r-2^i\) choices. On the other hand, each fixed \(d\)-dimensional subspace has
\[
        (2^d-1)(2^d-2)\cdots(2^d-2^{d-1})
\]
ordered bases, by the same argument inside \(\F_2^d\). Dividing gives the formula.
\end{proof}

\begin{lemma}[Uniform Gaussian-binomial estimate]\label{lem:gauss-estimate}
There is an absolute constant
\[
        C_2=\prod_{j=1}^{\infty}(1-2^{-j})^{-1}<\infty
\]
such that, for all \(0\leq d\leq r\),
\[
        2^{d(r-d)}\leq \gbinom{r}{d}\leq C_2\,2^{d(r-d)}.
\]
\end{lemma}

\begin{proof}
Starting from Lemma~\ref{lem:gauss-formula}, factor out the powers of \(2\):
\[
        \gbinom{r}{d}
        =2^{d(r-d)}
        \prod_{i=0}^{d-1}\frac{1-2^{i-r}}{1-2^{i-d}}.
\]
Since \(r\geq d\), one has \(2^{i-r}\leq 2^{i-d}\), hence
\[
        1-2^{i-r}\geq 1-2^{i-d}.
\]
Each factor in the product is therefore at least \(1\), proving the lower bound.

For the upper bound, use \(1-2^{i-r}\leq 1\). Then
\[
        \prod_{i=0}^{d-1}\frac{1-2^{i-r}}{1-2^{i-d}}
        \leq
        \prod_{i=0}^{d-1}(1-2^{i-d})^{-1}
        =
        \prod_{j=1}^{d}(1-2^{-j})^{-1}
        \leq
        \prod_{j=1}^{\infty}(1-2^{-j})^{-1}.
\]
The infinite product converges because \(\sum_{j\geq 1}2^{-j}<\infty\). This proves the estimate.
\end{proof}

\begin{lemma}[Number of subspaces of \(\F_2^r\)]\label{lem:subspace-total}
Let
\[
        G(r)=\sum_{d=0}^{r}\gbinom{r}{d}.
\]
Then
\[
        G(r)=2^{r^2/4+O(\log r)}.
\]
In particular,
\[
        \log_2 G(r)=\frac{r^2}{4}+O(\log r).
\]
\end{lemma}

\begin{proof}
The function \(d(r-d)\) is maximized at \(d=r/2\), with maximum \(r^2/4\). Therefore
Lemma~\ref{lem:gauss-estimate} gives
\[
        G(r)
        \leq
        \sum_{d=0}^r C_2 2^{d(r-d)}
        \leq
        (r+1)C_2 2^{r^2/4}.
\]
For the lower bound, take \(d=\lfloor r/2\rfloor\). Then
\[
        G(r)\geq \gbinom{r}{\lfloor r/2\rfloor}
        \geq 2^{\lfloor r/2\rfloor\lceil r/2\rceil}
        =2^{r^2/4+O(1)}.
\]
Combining the two estimates gives the claim.
\end{proof}

\begin{lemma}[Concentration of the dimension of a random subspace]\label{lem:dimension-concentration}
Let \(U_r\) be chosen uniformly from all linear subspaces of \(\F_2^r\). Then
\[
        \frac{\dim U_r}{r}\to \frac12
\]
in probability. Equivalently, for every fixed \(\varepsilon>0\),
\[
        \Prob\left(\left|\dim U_r-\frac r2\right|\geq \varepsilon r\right)\to 0.
\]
More quantitatively,
\[
        \Prob\left(\left|\dim U_r-\frac r2\right|\geq \varepsilon r\right)
        \leq
        (r+1)C_2\,2^{-\varepsilon^2r^2+O(1)}.
\]
\end{lemma}

\begin{proof}
For every \(d\),
\[
        d(r-d)=\frac{r^2}{4}-\left(d-\frac r2\right)^2.
\]
Thus, if \(|d-r/2|\geq \varepsilon r\), Lemma~\ref{lem:gauss-estimate} gives
\[
        \gbinom{r}{d}\leq C_2 2^{r^2/4-\varepsilon^2r^2}.
\]
There are at most \(r+1\) possible dimensions \(d\), so the number of exceptional subspaces is at
most
\[
        (r+1)C_2 2^{r^2/4-\varepsilon^2r^2}.
\]
By Lemma~\ref{lem:subspace-total}, the total number of subspaces is at least
\(2^{r^2/4+O(1)}\). Dividing the two bounds proves the quantitative estimate and hence the
probabilistic convergence.
\end{proof}

\section{The elementary lower bound for \texorpdfstring{\(s(n)\)}{s(n)}}

Let
\[
        m=\left\lfloor \frac n2\right\rfloor.
\]
Inside \(S_n\), define
\[
        E_m=\langle (1\,2),(3\,4),\ldots,(2m-1\,\,2m)\rangle.
\]
Then \(E_m\cong C_2^m\), because the displayed transpositions are disjoint and independent.

\begin{proposition}[Elementary lower bound]\label{prop:lower}
One has
\[
        s(n)\geq 2^{n^2/16+O(n)}.
\]
In fact,
\[
        s(n)\geq G\left(\left\lfloor \frac n2\right\rfloor\right),
\]
where \(G(r)\) is the number of subspaces of \(\F_2^r\).
\end{proposition}

\begin{proof}
The subgroup \(E_m\) is naturally an \(m\)-dimensional vector space over \(\F_2\): multiplication
in \(E_m\) corresponds to addition of exponent vectors in \(\F_2^m\). Therefore the subgroups of
\(E_m\) are exactly the \(\F_2\)-linear subspaces of \(\F_2^m\). Distinct subspaces give distinct
subgroups of \(S_n\). Hence
\[
        s(n)\geq |\Sub(E_m)|=G(m).
\]
By Lemma~\ref{lem:subspace-total},
\[
        G(m)=2^{m^2/4+O(\log m)}.
\]
Since \(m=\lfloor n/2\rfloor\), we have
\[
        \frac{m^2}{4}=\frac{n^2}{16}+O(n).
\]
Thus
\[
        s(n)\geq 2^{n^2/16+O(n)}.
\]
\end{proof}

\section{The upper bound and the asymptotic formula}

The next theorem is the non-elementary input, quoted from Roney-Dougal and Tracey \cite[Theorem 1]{RDT} rather than reproved.

\begin{importedtheorem}[Roney-Dougal--Tracey enumeration theorem]\label{thm:rdt-enum}
There exist absolute constants \(\alpha>0\) and \(\beta\) such that, for all integers \(n>1\),
\[
        2^{n^2/16+\alpha n\log n}
        \leq
        |\Sub(S_n)|
        \leq
        2^{n^2/16+\beta n^{3/2}}.
\]
\end{importedtheorem}

\begin{remark}
Pyber's 1993 work gave a weaker upper bound of the form \(2^{\xi n^2+o(n^2)}\), where
\(\xi=(\log_2 24)/6\), and conjectured the sharper constant \(1/16\). Theorem~\ref{thm:rdt-enum}
settles that conjecture. Thus the matching upper bound in this note should be attributed to
Roney-Dougal and Tracey, not to Pyber alone.
\end{remark}

\begin{proof}[Proof of Theorem~\ref{thm:main}]
The elementary lower bound from Proposition~\ref{prop:lower} gives
\[
        s(n)\geq 2^{n^2/16+O(n)}=2^{n^2/16+o(n^2)}.
\]
The upper bound in Theorem~\ref{thm:rdt-enum} gives
\[
        s(n)\leq 2^{n^2/16+\beta n^{3/2}}=2^{n^2/16+o(n^2)}.
\]
Combining the two inequalities yields
\[
        s(n)=2^{n^2/16+o(n^2)}.
\]
Taking logarithms gives
\[
        \log_2 s(n)=\frac{n^2}{16}+o(n^2),
\]
as claimed.
\end{proof}

\section{A global statistical theorem on subgroup order}\label{sec:global-order}

We now state the known global probabilistic input about uniformly random subgroups of \(S_n\), again quoted from Roney-Dougal and Tracey \cite[Theorem 6]{RDT}.

\begin{importedtheorem}[Roney-Dougal--Tracey random subgroup theorem]\label{thm:rdt-random}
Let \(H_n\) be uniformly random from \(\Sub(S_n)\). Let \(\mu\in[0,1/16)\) and let
\[
        \nu\in\left[0,\frac12-\frac{\sqrt3}{4}\right).
\]
Then, with probability tending to \(1\) as \(n\to\infty\), the group \(H_n\) has an elementary
abelian \(2\)-section of order at least \(2^{\mu n}\), and a Sylow \(2\)-subgroup of order at least
\(2^{\nu n}\).
\end{importedtheorem}

\begin{proof}[Proof of Theorem~\ref{thm:global-order}]
Let \(\varepsilon>0\). If \(\varepsilon\geq\rho\), the asserted lower bound is trivial because
\(|H_n|\geq 1\). If \(0<\varepsilon<\rho\), put \(\nu=\rho-\varepsilon\). Then
\(0\leq \nu<\rho\), so Theorem~\ref{thm:rdt-random} gives, with probability tending to \(1\), a
Sylow \(2\)-subgroup \(P_n\leq H_n\) satisfying
\[
        |P_n|\geq 2^{\nu n}=2^{(\rho-\varepsilon)n}.
\]
Since \(|P_n|\leq |H_n|\), it follows that
\[
        \Prob\left(\log_2 |H_n|\geq (\rho-\varepsilon)n\right)\to 1.
\]
The deterministic upper bound follows from Lagrange's theorem:
\[
        |H_n|\leq |S_n|=n!,
\]
and from the elementary inequality \(n!\leq n^n\), which gives \(\log_2(n!)\leq n\log_2 n\).
\end{proof}

\begin{remark}
The elementary abelian section conclusion in Theorem~\ref{thm:rdt-random} also implies
\(|H_n|\geq 2^{\mu n}\) with high probability for every \(\mu<1/16\), since the order of a section
of a finite group divides the order of the group. The Sylow \(2\)-subgroup assertion gives the
slightly larger explicit lower-tail constant \(\rho=1/2-\sqrt3/4\).
\end{remark}

\section{The order law inside the elementary abelian model}\label{sec:model-order}

The lower-bound construction not only gives many subgroups; it also gives a clean order statistic
inside that model.

\begin{theorem}[Model order theorem]\label{thm:model-order}
Let \(m=\lfloor n/2\rfloor\), and choose a subgroup \(U_n\leq E_m\) uniformly at random from all
subgroups of \(E_m\). Then
\[
        \log_2 |U_n|=\frac n4+o(n)
\]
in probability. Equivalently,
\[
        |U_n|=2^{n/4+o(n)}
\]
in probability.
\end{theorem}

\begin{proof}
Under the identification \(E_m\cong \F_2^m\), the subgroup \(U_n\) is a uniformly random subspace
of \(\F_2^m\). Its group order is
\[
        |U_n|=2^{\dim U_n},
\]
so
\[
        \log_2 |U_n|=\dim U_n.
\]
By Lemma~\ref{lem:dimension-concentration},
\[
        \dim U_n=\frac m2+o(m)
\]
in probability. Since \(m=n/2+O(1)\), this becomes
\[
        \log_2 |U_n|=\frac n4+o(n),
\]
which is the claim.
\end{proof}

\section{Why the global order distribution is a finer question}\label{sec:dihedral}

It is tempting to infer from Theorem~\ref{thm:model-order} that a uniformly random subgroup of
\(S_n\) itself should have order \(2^{n/4+o(n)}\). That inference is not justified by the leading
asymptotic formula. The following construction shows why: another natural family has the same
leading \(2^{n^2/16}\)-scale but a different linear order scale.

\begin{lemma}[A dihedral block group]\label{lem:dihedral-block}
Let \(D_8\) be the dihedral group of order \(8\), acting faithfully on four points as the symmetry
group of a square. Its Frattini subgroup \(\Phi(D_8)\), namely the intersection of all maximal
subgroups of \(D_8\), has order \(2\), and
\[
        D_8/\Phi(D_8)\cong C_2^2.
\]
\end{lemma}

\begin{proof}
Write
\[
        D_8=\langle r,s:r^4=s^2=1,\;srs=r^{-1}\rangle.
\]
The subgroup \(\langle r^2\rangle\) is central of order \(2\). The quotient
\(D_8/\langle r^2\rangle\) is generated by the images of \(r\) and \(s\), both of order \(2\), and
these images commute because the commutator relation becomes trivial modulo \(\langle r^2\rangle\).
Thus the quotient has order \(4\) and is isomorphic to \(C_2^2\).

It remains to identify \(\langle r^2\rangle\) with the Frattini subgroup. Since \(|D_8|=8\), the
maximal subgroups have order \(4\). They are exactly
\[
        \langle r\rangle,
        \qquad
        \langle r^2,s\rangle,
        \qquad
        \langle r^2,rs\rangle.
\]
Indeed, \(\langle r\rangle\) is the unique cyclic subgroup of order \(4\), while every subgroup of
order \(4\) not equal to \(\langle r\rangle\) contains a reflection and \(r^2\), giving one of the
two displayed Klein four subgroups according to whether the reflection is \(r^{2j}s\) or
\(r^{2j+1}s\). Their intersection is precisely \(\langle r^2\rangle\). Hence
\[
        \Phi(D_8)=\langle r\rangle\cap \langle r^2,s\rangle\cap\langle r^2,rs\rangle
        =\langle r^2\rangle.
\]
The quotient by \(\Phi(D_8)\) is therefore \(C_2^2\), as proved above.
\end{proof}

\begin{proposition}[Many subgroups at order scale \(2^{n/2}\)]\label{prop:dihedral-family}
Suppose \(n=4k\). Then \(S_n\) contains a family \(\mathcal F_n\) of subgroups with
\[
        |\mathcal F_n|=2^{n^2/16+O(\log n)}
\]
such that all but an exponentially small proportion of the members of \(\mathcal F_n\) satisfy
\[
        \log_2 |H|=\frac n2+o(n).
\]
\end{proposition}

\begin{proof}
Partition \(\{1,\ldots,n\}\) into \(k\) fixed blocks of size \(4\). On each block let a copy of
\(D_8\) act as in Lemma~\ref{lem:dihedral-block}, and set
\[
        P_k=D_8^k\leq S_{4k}.
\]
Let
\[
        N_k=\Phi(D_8)^k\trianglelefteq P_k.
\]
Then
\[
        N_k\cong C_2^k
        \qquad\text{and}\qquad
        P_k/N_k\cong C_2^{2k}.
\]
Let
\[
        \pi:P_k\to P_k/N_k
\]
be the quotient map. For every subspace \(U\leq P_k/N_k\), the inverse image
\[
        H_U=\pi^{-1}(U)
\]
is a subgroup of \(P_k\), hence a subgroup of \(S_{4k}\). Distinct subspaces \(U\) give distinct
subgroups \(H_U\), since \(\pi(H_U)=U\). Therefore the number of such subgroups is exactly the
number of subspaces of \(\F_2^{2k}\), namely
\[
        G(2k)=2^{(2k)^2/4+O(\log k)}=2^{k^2+O(\log k)}.
\]
Since \(n=4k\), this is
\[
        2^{n^2/16+O(\log n)}.
\]

Moreover, if \(\dim U=d\), then
\[
        |H_U|=|N_k|\,|U|=2^k 2^d=2^{k+d}.
\]
By Lemma~\ref{lem:dimension-concentration}, a uniformly random subspace \(U\leq\F_2^{2k}\) has
\[
        d=k+o(k)
\]
in probability. Hence, for all but an exponentially small proportion of the subspaces \(U\),
\[
        \log_2 |H_U|=k+d=2k+o(k)=\frac n2+o(n).
\]
This proves the claim.
\end{proof}

\begin{corollary}[The leading formula alone does not give an exact global order law]\label{cor:order-caution}
The asymptotic formula
\[
        s(n)=2^{n^2/16+o(n^2)}
\]
does not by itself imply that a uniformly random subgroup of \(S_n\) has order
\(2^{n/4+o(n)}\). Indeed, for \(4\mid n\), there is a family of
\[
        2^{n^2/16+O(\log n)}
\]
subgroups of order scale \(2^{n/2+o(n)}\), as shown in Proposition~\ref{prop:dihedral-family}.
\end{corollary}

\begin{proof}
Theorem~\ref{thm:model-order} proves the order scale \(2^{n/4+o(n)}\) only for subgroups chosen
inside a fixed elementary abelian subgroup of rank \(\lfloor n/2\rfloor\). Proposition~\ref{prop:dihedral-family}
constructs, on the same leading \(2^{n^2/16}\) logarithmic scale, subgroups whose order is instead
\(2^{n/2+o(n)}\). Therefore the leading exponential count cannot distinguish the global
probability weights of these different families. A sharper theorem for a uniformly random subgroup
of \(S_n\) requires finer enumeration than the main logarithmic asymptotic.
\end{proof}

\section{A useful rank check for elementary abelian 2-subgroups}

The following elementary observation explains why the rank \(\lfloor n/2\rfloor\) appearing in the
lower bound is the correct maximum possible rank for elementary abelian \(2\)-subgroups of \(S_n\).

\begin{lemma}\label{lem:rank-bound}
If \(A\leq S_n\) is elementary abelian of exponent \(2\), say \(A\cong C_2^r\), then
\[
        r\leq \left\lfloor\frac n2\right\rfloor.
\]
\end{lemma}

\begin{proof}
Let \(A\) act on \(\{1,\ldots,n\}\), and decompose the set into \(A\)-orbits
\[
        \Omega_1,\ldots,\Omega_t.
\]
For each orbit \(\Omega_i\), the image of \(A\) in \(\Sym(\Omega_i)\) is a transitive abelian
permutation group. A transitive abelian permutation group acts regularly after factoring out its
kernel: if an element fixes one point of a transitive orbit, then, because the group is abelian, it
fixes every point of that orbit. Hence the faithful image on \(\Omega_i\) is an elementary abelian
2-group of order \(|\Omega_i|\). Thus \(|\Omega_i|=2^{a_i}\) for some \(a_i\geq 0\), and the rank
of the image on \(\Omega_i\) is \(a_i\).

The diagonal action of \(A\) on all its orbits is faithful, so the natural homomorphism from \(A\)
to the product of its images on the orbits is injective. Therefore
\[
        r\leq \sum_{i=1}^t a_i.
\]
For each orbit, \(a_i\leq 2^{a_i}/2\) when \(a_i\geq 1\), and the inequality is also true for
\(a_i=0\). Hence
\[
        r\leq \sum_{i=1}^t a_i
        \leq
        \frac12\sum_{i=1}^t 2^{a_i}
        =\frac n2.
\]
Since \(r\) is an integer, \(r\leq\lfloor n/2\rfloor\).
\end{proof}

\section{Verification of the answer}

We finish by isolating exactly what has been proved and where each input enters.

\begin{enumerate}[label=\textup{(\arabic*)},leftmargin=2.5em]
    \item Lemmas~\ref{lem:gauss-formula}--\ref{lem:dimension-concentration} are fully elementary
    linear algebra over \(\F_2\). They prove both the count and dimension concentration for
    subspaces of \(\F_2^r\).

    \item Proposition~\ref{prop:lower} injects the subgroup lattice of
    \(C_2^{\lfloor n/2\rfloor}\) into \(\Sub(S_n)\). This proves
    \[
            s(n)\geq 2^{n^2/16+O(n)}.
    \]

    \item The non-elementary upper bound is Theorem~\ref{thm:rdt-enum}, due to Roney-Dougal and
    Tracey. Combining it with the elementary lower bound gives
    \[
            s(n)=2^{n^2/16+o(n^2)}.
    \]
    This is the requested asymptotic formula.

    \item The genuine global statistical theorem on order is Theorem~\ref{thm:global-order}: if
    \(H_n\) is uniformly random in \(\Sub(S_n)\), then for every \(\varepsilon>0\),
    \[
            \Prob\left(\log_2|H_n|\geq
            \left(\frac12-\frac{\sqrt3}{4}-\varepsilon\right)n\right)\to 1.
    \]
    This is obtained from the Roney-Dougal--Tracey random subgroup theorem.

    \item Theorem~\ref{thm:model-order} is a fully proved sharper statistical theorem for the
    elementary abelian lower-bound model. It says that a random subgroup of
    \(C_2^{\lfloor n/2\rfloor}\leq S_n\) has order \(2^{n/4+o(n)}\).

    \item Proposition~\ref{prop:dihedral-family} and Corollary~\ref{cor:order-caution} explain why
    the leading logarithmic enumeration alone does not force a unique exact global order scale.
\end{enumerate}

Thus the original question is answered as follows: the number of subgroups of \(S_n\) is
\(2^{n^2/16+o(n^2)}\); there is a global high-probability theorem proving that a uniformly random
subgroup has exponentially large order, and there is an exact order law in the elementary abelian
model that supplies the elementary lower bound. A sharper global concentration theorem for
\(\log_2|H_n|\), if desired, requires finer information than the leading asymptotic count.

\begin{thebibliography}{9}

\bibitem{Pyber}
L. Pyber,
\emph{Enumerating finite groups of given order},
Ann. of Math. (2) \textbf{137} (1993), 203--220.

\bibitem{RDT}
C. M. Roney-Dougal and G. Tracey,
\emph{Subgroups of symmetric groups: enumeration and asymptotic properties},
arXiv:2503.05416, 2025.

\end{thebibliography}

\end{document}
